% ASSERT_CONVENTION: natural_units=natural, metric_signature=mostly_minus, fourier_convention=physics, coupling_convention=alpha_s, generator_normalization=T(fund)=1/2, covariant_derivative_sign=D_mu=partial_mu-igA
%
% Paper 4: Why Three Generations -- The n_g >= 3 Bound Across Duality Types
%
% Conventions (Seiberg hep-th/9411149 + Intriligator-Seiberg hep-th/9509066):
%   T(fund SU(N)) = 1/2 per Weyl fermion
%   T(vector SO(N)) = 1
%   T(fund Sp(N)) = 1/2
%   Sp(N_c) = rank N_c, fund dim 2N_c.  Sp(1) = SU(2).
%   h(SU(N)) = N, h(SO(N)) = N-2, h(Sp(N)) = N+1
%   b_0(SU(N_c)) = (11N_c - 2N_f)/3 for N_f Dirac flavours (no adjoint)
%   b_0(SO(N_c)) = 3(N_c - 2) - N_f for N_f vectors
%   b_0(Sp(N_c)) = 3(N_c + 1) - N_f for 2N_f fundamentals
%   N_f counts Dirac flavour pairs for SU; vectors for SO; 2N_f = total funds for Sp
%   Non-SUSY extrapolation: ALL claims carry conjectural qualifiers
%
\documentclass[11pt,a4paper]{article}

%% ---------- Packages ----------
\usepackage{amsmath,amssymb,amsthm,mathtools}
\usepackage{hyperref}
\usepackage[capitalise,noabbrev]{cleveref}
\usepackage{enumitem}
\usepackage{booktabs}
\usepackage{array}
\usepackage{caption}
\usepackage{xcolor}
\usepackage{slashed}
\usepackage{cite}
\usepackage[margin=2.5cm]{geometry}

%% ---------- Hyperref configuration ----------
\hypersetup{
  colorlinks=true,
  linkcolor=blue!70!black,
  citecolor=green!50!black,
  urlcolor=blue!80!black,
}

%% ---------- Theorem environments ----------
\theoremstyle{plain}
\newtheorem{theorem}{Theorem}[section]
\newtheorem{lemma}[theorem]{Lemma}
\newtheorem{proposition}[theorem]{Proposition}
\newtheorem{corollary}[theorem]{Corollary}
\theoremstyle{definition}
\newtheorem{definition}[theorem]{Definition}
\theoremstyle{remark}
\newtheorem{remark}[theorem]{Remark}

%% ---------- Physics macros ----------
\newcommand{\Nc}{N_{\!c}}
\newcommand{\Nf}{N_{\!f}}
\newcommand{\ngen}{n_g}
\newcommand{\SU}[1]{\mathrm{SU}(#1)}
\newcommand{\UU}[1]{\mathrm{U}(#1)}
\newcommand{\SO}[1]{\mathrm{SO}(#1)}
\newcommand{\Sp}[1]{\mathrm{Sp}(#1)}
\newcommand{\fund}{\square}
\newcommand{\antifund}{\overline{\square}}
\newcommand{\adj}{\mathrm{adj}}
\newcommand{\antisym}{\wedge^{\!2}}
\newcommand{\sym}{\mathrm{S}_2}
\newcommand{\Tr}{\mathrm{Tr}}
\newcommand{\ie}{\textit{i.e.}}
\newcommand{\eg}{\textit{e.g.}}
\newcommand{\cf}{\textit{cf.}}
\newcommand{\IR}{\ensuremath{\mathrm{IR}}}
\newcommand{\UV}{\ensuremath{\mathrm{UV}}}
\newcommand{\SM}{\ensuremath{\mathrm{SM}}}

%% ---------- Layout ----------
\setlength{\parskip}{0.4em}
\setlength{\parindent}{0em}

%% ================================================================
%%  DOCUMENT
%% ================================================================
\begin{document}

\title{\textbf{Why Three Generations:\\
The $\ngen \geq 3$ Bound Across Duality Types}}
\author{Notes from the Dualised Standard Model Programme}
\date{March 2026}
\maketitle

\begin{abstract}
The Cacciapaglia--Sannino construction identifies the Standard Model as the
magnetic (infrared) dual of an asymptotically free electric gauge theory,
yielding a lower bound $\ngen \geq 3$ on the number of fermion generations
from the requirement that the electric gauge group be non-abelian.  We
systematically derive the analogous generation bound for \emph{every}
established $\mathcal{N}=1$ Seiberg-type duality: $\SU{\Nc}$ (Seiberg),
$\SO{\Nc}$ and $\Sp{\Nc}$ (Intriligator--Seiberg), and $\SU{\Nc}$ with
adjoint matter (Kutasov--Schwimmer).  For each duality type, we identify
the electric gauge group as a function of~$\ngen$, impose the non-abelian
condition, and check compatibility with the conformal window.  Our main
results are: (i)~the $\SU{}$ and $\Sp{}$ routes independently constrain
$\ngen \geq 3$ and $2 \leq \ngen \leq 3$ respectively, and their intersection
\emph{uniquely} selects $\ngen = 3$; (ii)~the $\SO{}$ route requires
$\ngen \geq 4$ for conformal window compatibility, excluding it as a
duality frame for the observed Standard Model; (iii)~Kutasov--Schwimmer
duality with $k \geq 2$ admits no integer generation number in the
conformal window, ruling out higher-$k$ adjoint deformations.  Within the chosen duality-embedding framework, and conditional on importing SUSY conformal-window formulas, $\ngen = 3$ emerges as the unique
solution compatible with the intersection of all viable Seiberg-type
dualities.  All claims regarding non-supersymmetric duality are
conjectural.
\end{abstract}

\tableofcontents
\newpage

%% ================================================================
%%  SECTION 1: INTRODUCTION
%% ================================================================
\section{Introduction}
\label{sec:introduction}

The Standard Model of particle physics contains three generations of
quarks and leptons.  No compelling explanation for this number exists
within the Standard Model itself; the generation structure is an input
rather than a prediction.  The question ``why three generations?'' has
been called one of the deepest puzzles in fundamental
physics~\cite{Frampton1992}.

A novel perspective on this question emerged from the programme of
Cacciapaglia and Sannino~\cite{CacciapagliaSannino2024,Sannino2011},
who proposed that the entire Standard Model can be understood as the
\emph{magnetic} (infrared) dual of an asymptotically free electric gauge
theory, extrapolating the pattern of Seiberg
duality~\cite{Seiberg1995} beyond supersymmetry.  In their construction,
the QCD gauge group $\SU{3}_c$ is the magnetic dual of an $\SU{N}$
electric theory embedded in a Pati--Salam~\cite{PS1974} framework, with
the electric gauge group determined as $\SU{2\ngen - 4}$.  The requirement
that this group be non-abelian ($2\ngen - 4 \geq 2$) immediately yields
\begin{equation}\label{eq:CS-bound}
  \ngen \geq 3\,.
\end{equation}
If the conjectured non-supersymmetric duality
holds~\cite{AMPS2011,Sannino2011}, the existence of at least three
generations is not an accident but a mathematical necessity for the
electric theory to be a non-trivial gauge theory.

This is a striking result, but it relies on a specific choice of duality
type ($\SU{\Nc}$ Seiberg duality in the Pati--Salam embedding).  The
landscape of established $\mathcal{N}=1$ Seiberg-type dualities is far
richer, including:
\begin{itemize}
  \item $\SU{\Nc}$ with fundamentals (Seiberg~\cite{Seiberg1995}),
  \item $\SO{\Nc}$ with vectors
    (Intriligator--Seiberg~\cite{IS9506101}),
  \item $\Sp{\Nc}$ with fundamentals
    (Intriligator--Seiberg~\cite{IS9506101}),
  \item $\SU{\Nc}$ with fundamentals and adjoint
    (Kutasov--Schwimmer~\cite{KS9505004}).
\end{itemize}
Each duality type has a different map between electric and magnetic gauge
groups, different conformal windows, and potentially different generation
bounds.

In this paper, we systematically derive the generation bound for each of
these duality types.  The logic is uniform: identify a Standard Model
gauge group sector as the magnetic theory, use the duality map to
determine the electric gauge group as a function of~$\ngen$, require the
electric group to be non-abelian, and check conformal window
compatibility.  We then ask: \emph{does the intersection of all bounds
uniquely select $\ngen = 3$?}

Previous work~\cite{Paper1} established that the intersection of the
$\SU{}$ (Cacciapaglia--Sannino) and $\Sp{}$ (Intriligator--Seiberg)
routes selects $\ngen = 3$ uniquely, using the fact that the Sp~route
provides an \emph{upper} bound $\ngen \leq 3$ from conformal window
constraints.  Here we extend this analysis to include the $\SO{}$ and
Kutasov--Schwimmer dualities, providing a complete survey across all
classical group duality types.

\medskip

\textbf{Epistemic convention.}  The SUSY Seiberg-type dualities are
established results of $\mathcal{N}=1$ supersymmetric gauge theory.  All
claims extrapolating these dualities to the non-supersymmetric Standard
Model are \textbf{conjectural} and carry an implicit qualifier ``if the
non-SUSY duality holds.''  The generation bounds derived here are exact
consequences of the group-theoretic structure; the physical relevance
to the Standard Model depends on the dynamical assumption that such
dualities extend beyond SUSY.

\medskip

\textbf{Summary of results.}
\begin{enumerate}
  \item The $\SU{}$ Seiberg route and the $\SU{}$ Pati--Salam route both
    give $\ngen \geq 3$, with conformal window compatibility for
    $\ngen \in \{3, 4\}$ and $\ngen \in \{3, 4, 5, 6\}$ respectively
    (\S\ref{sec:SU-duality}).
  \item The $\Sp{}$ route gives $\ngen \geq 2$ from group theory and
    $\ngen \leq 3$ from the conformal window, yielding
    $2 \leq \ngen \leq 3$
    (\S\ref{sec:Sp-duality}).
  \item The $\SO{}$ route gives $\ngen \geq 3$ from group theory but
    requires $\ngen \geq 4$ for conformal window compatibility,
    \emph{excluding} $\ngen = 3$
    (\S\ref{sec:SO-duality}).
  \item Kutasov--Schwimmer duality with $k \geq 2$ admits no integer
    $\ngen$ in the conformal window; only $k = 1$ (recovering Seiberg
    duality) is viable (\S\ref{sec:KS-duality}).
  \item The intersection of all viable duality-type bounds uniquely
    selects $\ngen = 3$ (\S\ref{sec:master-table}).
\end{enumerate}


%% ================================================================
%%  SECTION 2: FRAMEWORK AND CONVENTIONS
%% ================================================================
\section{Framework and Conventions}
\label{sec:framework}

\subsection{The duality programme}

We follow the framework of Cacciapaglia and
Sannino~\cite{CacciapagliaSannino2024}: a Standard Model gauge group
factor $G_{\SM}$ is identified as the \emph{magnetic} gauge group of a
Seiberg-type duality, and the electric theory provides the UV completion.
The matter content of the magnetic theory is fixed by the Standard Model
spectrum, which determines the number of flavours~$\Nf$ as a function
of the generation number~$\ngen$.

The generation bound arises from two independent constraints:
\begin{enumerate}
  \item \textbf{Non-abelian condition:} The electric gauge group must be
    non-abelian (\ie, non-trivial as a gauge theory).  For $\SU{N}$,
    this means $N \geq 2$; for $\SO{N}$, $N \geq 3$; for $\Sp{N}$,
    $N \geq 1$ (where $\Sp{N}$ has rank~$N$).
  \item \textbf{Conformal window:} Both the electric and magnetic
    theories must lie within or at the boundary of the conformal window
    for the duality to produce a non-trivial interacting fixed point.
\end{enumerate}

\subsection{Conventions}

We follow the conventions of Seiberg~\cite{Seiberg1995} and
Intriligator--Seiberg~\cite{IS9506101,IS1995}:
\begin{itemize}
  \item $\Sp{\Nc}$ denotes the symplectic group of rank~$\Nc$, with
    fundamental representation of dimension~$2\Nc$.  Thus
    $\Sp{1} = \SU{2}$.
  \item $\SO{\Nc}$ has an $\Nc$-dimensional vector representation.
  \item Dynkin indices: $T(\mathrm{fund}\;\SU{N}) = \tfrac{1}{2}$,
    $T(\mathrm{vector}\;\SO{N}) = 1$,
    $T(\mathrm{fund}\;\Sp{N}) = \tfrac{1}{2}$.
  \item Dual Coxeter numbers: $h(\SU{N}) = N$, $h(\SO{N}) = N - 2$,
    $h(\Sp{N}) = N + 1$.
  \item $\Nf$ counts Dirac flavour pairs for SU duality, vector fields
    for SO duality, and $2\Nf$ total fundamentals for Sp duality.
\end{itemize}

\subsection{SM matter content as a function of $\ngen$}
\label{ssec:SM-matter}

The Standard Model with $\ngen$ generations contains:
\begin{itemize}
  \item \textbf{QCD sector ($\SU{3}_c$):} $2\ngen$ Dirac quarks
    (each generation contributes one up-type and one down-type quark),
    hence $\Nf^{\mathrm{QCD}} = 2\ngen$.
  \item \textbf{EW sector ($\SU{2}_L$):} $4\ngen$ Weyl doublets per
    generation (3~quark doublets from colour $+$ 1~lepton doublet),
    hence $2\Nf^{\mathrm{EW}} = 4\ngen$, giving
    $\Nf^{\mathrm{EW}} = 2\ngen$ in the Sp counting convention.
  \item \textbf{Pati--Salam ($\SU{4}_{\mathrm{PS}}$):}
    $2\ngen$ Dirac ``quarks'' in the fundamental of $\SU{4}$, but the
    specific embedding introduces an additional offset (see
    \S\ref{ssec:pati-salam}).
\end{itemize}


%% ================================================================
%%  SECTION 3: SU(N_c) SEIBERG DUALITY
%% ================================================================
\section{$\SU{\Nc}$ Seiberg Duality}
\label{sec:SU-duality}

\subsection{Seiberg duality map}

The electric theory is $\mathcal{N}=1$ $\SU{\Nc}$ SQCD with $\Nf$
fundamental flavours.  The magnetic dual is $\SU{\Nf - \Nc}$ with $\Nf$
dual quarks and an $\Nf \times \Nf$ meson
matrix~\cite{Seiberg1995}.

The conformal window is
\begin{equation}\label{eq:SU-CW}
  \frac{3}{2}\,\Nc < \Nf < 3\Nc\,.
\end{equation}

\subsection{Direct QCD route}
\label{ssec:direct-QCD}

We identify $\SU{3}_c$ as the magnetic gauge group:
\begin{equation}\label{eq:SU-magnetic}
  \Nf - \Nc = 3
  \qquad\Longrightarrow\qquad
  \Nc = \Nf - 3\,.
\end{equation}
The number of QCD Dirac flavours is $\Nf = 2\ngen$
(\cf~\S\ref{ssec:SM-matter}).  The electric gauge group is therefore
\begin{equation}\label{eq:SU-electric}
  \boxed{\SU{2\ngen - 3}}\,.
\end{equation}

\paragraph{Non-abelian condition.}
$2\ngen - 3 \geq 2$ requires
\begin{equation}\label{eq:SU-bound}
  \boxed{\ngen \geq \frac{5}{2}
  \qquad\Longrightarrow\qquad \ngen \geq 3}\,.
\end{equation}

\paragraph{Conformal window.}
From~\eqref{eq:SU-CW} with $\Nc = 2\ngen - 3$:
\begin{align}
  \text{Lower:}&\quad \frac{3(2\ngen - 3)}{2} < 2\ngen
  \quad\Longrightarrow\quad
  6\ngen - 9 < 4\ngen
  \quad\Longrightarrow\quad
  \ngen < \frac{9}{2}\,,
  \label{eq:SU-CW-lower} \\
  \text{Upper:}&\quad 2\ngen < 3(2\ngen - 3) = 6\ngen - 9
  \quad\Longrightarrow\quad
  \ngen > \frac{9}{4}\,.
  \label{eq:SU-CW-upper}
\end{align}
So the conformal window requires $\frac{9}{4} < \ngen < \frac{9}{2}$,
\ie,
\begin{equation}\label{eq:SU-CW-ng}
  \ngen \in \{3, 4\}\,.
\end{equation}

\paragraph{Numerical check ($\ngen = 3$).}
$\Nc = 3$, $\Nf = 6$.  Electric $\SU{3}$ with 6~flavours.
Conformal window: $4.5 < 6 < 9$.  \checkmark

\subsection{Pati--Salam route (Cacciapaglia--Sannino)}
\label{ssec:pati-salam}

In the Pati--Salam embedding~\cite{PS1974,CacciapagliaSannino2024}, the
electric theory is an $\SU{N}$ gauge theory with $\Nf$ Dirac
fundamentals \emph{and} one adjoint Weyl fermion~$\lambda$.  The
magnetic theory contains $\SU{3}_c$ as a subgroup of $\SU{4}_{\mathrm{PS}}$,
with the ``fourth colour'' corresponding to lepton number.

The specific Pati--Salam embedding yields an electric gauge group of
\begin{equation}\label{eq:PS-electric}
  \boxed{\SU{2\ngen - 4}}
\end{equation}
(see~\cite{CacciapagliaSannino2024} and Lecture~7 of the accompanying
notes).

\paragraph{Non-abelian condition.}
$2\ngen - 4 \geq 2$ requires
\begin{equation}\label{eq:PS-bound}
  \boxed{\ngen \geq 3}\,.
\end{equation}

\paragraph{Conformal window.}
The beta function with an adjoint is
$b_0 = 2\Nc - \Nf$ (using $3\Nc - \Nc = 2\Nc$ from the adjoint
contribution).  The conformal window for the KS duality at $k = 1$
(\ie, massive adjoint, reducing to Seiberg) is
$\frac{3}{2}\Nc < \Nf < 3\Nc$, identical to~\eqref{eq:SU-CW}.
With $\Nc = 2\ngen - 4$ and $\Nf = 2\ngen$:
\begin{align}
  \text{Lower:}&\quad 3(2\ngen - 4)/2 < 2\ngen
  \;\Longrightarrow\;
  3\ngen - 6 < 2\ngen
  \;\Longrightarrow\;
  \ngen < 6\,,
  \label{eq:PS-CW-lower} \\
  \text{Upper:}&\quad 2\ngen < 3(2\ngen - 4) = 6\ngen - 12
  \;\Longrightarrow\;
  \ngen > 3\,.
  \label{eq:PS-CW-upper}
\end{align}
Thus $3 < \ngen < 6$, giving $\ngen \in \{4, 5\}$ strictly inside the
conformal window.

For $\ngen = 3$: $\Nc = 2$, $\Nf = 6$.  Lower boundary:
$\frac{3}{2} \times 2 = 3 < 6$ (\checkmark); upper: $6 < 6$
(\emph{not} satisfied).  So $\ngen = 3$ is at the upper boundary of
asymptotic freedom, corresponding to the free electric phase.  In
Sannino's analysis~\cite{Sannino2011}, this boundary case is still
considered viable because the magnetic theory ($\SU{3}_c$) is
asymptotically free and the duality relation connects the two theories
in the vicinity of the conformal window.

Combined PS constraint (non-abelian + approximate CW):
\begin{equation}\label{eq:PS-combined}
  3 \leq \ngen \leq 6\,.
\end{equation}


%% ================================================================
%%  SECTION 4: Sp(N_c) DUALITY
%% ================================================================
\section{$\Sp{\Nc}$ Duality}
\label{sec:Sp-duality}

\subsection{Sp duality map}

The electric theory is $\mathcal{N}=1$ $\Sp{\Nc}$ with $2\Nf$
fundamentals~\cite{IS9506101}.  The magnetic dual is $\Sp{\Nf - \Nc - 2}$
with $2\Nf$ dual fundamentals and an antisymmetric meson.

The conformal window is
\begin{equation}\label{eq:Sp-CW}
  \frac{3(\Nc + 1)}{2} < \Nf < 3(\Nc + 1)\,.
\end{equation}

\subsection{EW sector: $\SU{2}_L = \Sp{1}$ as magnetic}

We identify $\SU{2}_L \cong \Sp{1}$ as the magnetic gauge group
(following~\cite{Paper1}):
\begin{equation}\label{eq:Sp-constraint}
  \Nf - \Nc - 2 = 1
  \qquad\Longrightarrow\qquad
  \Nc = \Nf - 3\,.
\end{equation}

The electroweak sector contains $2\Nf = 4\ngen$ Weyl doublets (see
\S\ref{ssec:SM-matter}), giving $\Nf = 2\ngen$.
The electric gauge group is
\begin{equation}\label{eq:Sp-electric}
  \boxed{\Sp{2\ngen - 3}}\,.
\end{equation}

\paragraph{Non-abelian condition.}
$\Sp{\Nc}$ requires $\Nc \geq 1$, so $2\ngen - 3 \geq 1$:
\begin{equation}\label{eq:Sp-bound}
  \boxed{\ngen \geq 2}\,.
\end{equation}

\paragraph{Conformal window.}
From~\eqref{eq:Sp-CW} with $\Nc = 2\ngen - 3$:
\begin{align}
  \text{Lower:}&\quad \frac{3(2\ngen - 2)}{2} < 2\ngen
  \;\Longrightarrow\;
  3\ngen - 3 < 2\ngen
  \;\Longrightarrow\;
  \ngen < 3\,,
  \label{eq:Sp-CW-lower} \\
  \text{Upper:}&\quad 2\ngen < 3(2\ngen - 2) = 6\ngen - 6
  \;\Longrightarrow\;
  \ngen > \frac{3}{2}\,.
  \label{eq:Sp-CW-upper}
\end{align}
So $\frac{3}{2} < \ngen < 3$ inside the conformal window, giving
\begin{equation}
  \ngen = 2 \quad\text{(inside CW)}\,.
\end{equation}

For $\ngen = 3$: $\Nc = 3$, $\Nf = 6$.  Lower boundary:
$\frac{3}{2}(3 + 1) = 6 = \Nf$.  So $\ngen = 3$ sits \emph{exactly at
the lower boundary} of the conformal window (the free magnetic phase
boundary), where the magnetic description is weakly coupled and the
duality is at its most useful~\cite{Paper1}.

The combined Sp constraint is therefore:
\begin{equation}\label{eq:Sp-combined}
  \boxed{2 \leq \ngen \leq 3}\,.
\end{equation}


%% ================================================================
%%  SECTION 5: SO(N_c) DUALITY
%% ================================================================
\section{$\SO{\Nc}$ Duality}
\label{sec:SO-duality}

\subsection{SO duality map}

The electric theory is $\mathcal{N}=1$ $\SO{\Nc}$ with $\Nf$
vectors~\cite{IS9506101}.  The magnetic dual is
\begin{equation}\label{eq:SO-dual}
  \SO{\Nf - \Nc + 4}
\end{equation}
with $\Nf$ dual vectors and a symmetric meson.  Note the crucial
$+4$ shift, which distinguishes SO from SU duality.

The conformal window is
\begin{equation}\label{eq:SO-CW}
  \frac{3(\Nc - 2)}{2} < \Nf < 3(\Nc - 2)\,.
\end{equation}

\subsection{Embedding: $\SO{6} \cong \SU{4}_{\mathrm{PS}}$ as magnetic}

The Standard Model gauge groups are $\SU{3}_c \times \SU{2}_L
\times \UU{1}_Y$.  None of these is an SO group.  However, in the
Pati--Salam framework, $\SU{3}_c$ extends to $\SU{4}_{\mathrm{PS}}$,
and we may use the local isomorphism
$\SO{6} \cong \SU{4}$ to embed the Pati--Salam sector in an SO
duality~\cite{IS9506101,PS1974}.

The 6-dimensional vector of $\SO{6}$ decomposes under
$\SU{3}_c \subset \SU{4}$ as $\mathbf{3} \oplus \overline{\mathbf{3}}$.
Each vector therefore accommodates one quark plus one antiquark.  With
$2\ngen$ Dirac quarks in the Standard Model, we need
\begin{equation}
  \Nf = 2\ngen
\end{equation}
vectors of $\SO{6}$.

Setting the magnetic gauge group to $\SO{6}$:
\begin{equation}
  \Nf - \Nc + 4 = 6
  \qquad\Longrightarrow\qquad
  \Nc = \Nf - 2 = 2\ngen - 2\,.
\end{equation}
The electric gauge group is
\begin{equation}\label{eq:SO-electric}
  \boxed{\SO{2\ngen - 2}}\,.
\end{equation}

\paragraph{Non-abelian condition.}
$\SO{\Nc}$ is non-abelian for $\Nc \geq 3$:
\begin{equation}\label{eq:SO-bound}
  2\ngen - 2 \geq 3
  \qquad\Longrightarrow\qquad
  \boxed{\ngen \geq \frac{5}{2} \;\Longrightarrow\; \ngen \geq 3}\,.
\end{equation}

\paragraph{Conformal window.}
From~\eqref{eq:SO-CW} with $\Nc = 2\ngen - 2$:
\begin{align}
  \text{Lower:}&\quad \frac{3(2\ngen - 4)}{2} < 2\ngen
  \;\Longrightarrow\;
  3\ngen - 6 < 2\ngen
  \;\Longrightarrow\;
  \ngen < 6\,,
  \label{eq:SO-CW-lower} \\
  \text{Upper:}&\quad 2\ngen < 3(2\ngen - 4) = 6\ngen - 12
  \;\Longrightarrow\;
  \ngen > 3\,.
  \label{eq:SO-CW-upper}
\end{align}
So the conformal window requires $3 < \ngen < 6$:
\begin{equation}\label{eq:SO-CW-ng}
  \ngen \in \{4, 5\} \quad\text{(strictly inside CW)}\,.
\end{equation}

\paragraph{The $\ngen = 3$ case.}
For $\ngen = 3$: $\Nc = 4$, $\Nf = 6$.  The one-loop beta function:
\begin{equation}
  b_0 = 3(\Nc - 2) - \Nf = 3(2) - 6 = 0\,.
\end{equation}
The electric $\SO{4}$ theory with 6~vectors has a \emph{vanishing}
one-loop beta function and is \emph{not} asymptotically free.  Moreover,
$\SO{4} \cong \SU{2} \times \SU{2}$ is not simple; it is a product of
two abelian-looking SU(2) factors.

This means the SO~route is \textbf{incompatible with $\ngen = 3$}: the
electric theory is neither asymptotically free nor inside the conformal
window.

\paragraph{Combined SO constraint.}
The non-abelian condition gives $\ngen \geq 3$, but the conformal window
requires $\ngen \geq 4$:
\begin{equation}\label{eq:SO-combined}
  \boxed{\ngen \geq 4 \quad\text{(SO route)}}\,.
\end{equation}
This is a stronger bound than SU or Sp, and it \emph{excludes} the
observed value $\ngen = 3$.  The SO~duality therefore cannot serve as the
duality frame for the Standard Model with three generations.

\begin{remark}[Physical interpretation]
\label{rem:SO-exclusion}
The exclusion of the SO~route has a clear physical origin: the shift
$+4$ in the SO duality map $\SO{\Nf - \Nc + 4}$ (compared to $+0$ for
SU) means the electric gauge group is \emph{smaller} relative to the
magnetic one than in the SU case.  For a given number of generations,
the electric SO group has too few ``colours'' relative to ``flavours''
to maintain asymptotic freedom.  This is a direct consequence of the
different group-theoretic structure of orthogonal vs.\ unitary groups.
\end{remark}


%% ================================================================
%%  SECTION 6: KUTASOV-SCHWIMMER DUALITY
%% ================================================================
\section{Kutasov--Schwimmer Duality}
\label{sec:KS-duality}

\subsection{KS duality map}

The electric theory is $\mathcal{N}=1$ $\SU{\Nc}$ with $\Nf$
fundamentals, one adjoint~$\Phi$, and superpotential
$W = \frac{g}{k+1}\Tr(\Phi^{k+1})$, where $k \geq 1$ is a positive
integer~\cite{KS9505004,KSS9507040}.

The magnetic dual is $\SU{k\Nf - \Nc}$ with $\Nf$ dual quarks, one
dual adjoint~$\phi$, and $k$~meson species $M_j \sim Q\Phi^j
\widetilde{Q}$ ($j = 0, \ldots, k-1$).

The conformal window is~\cite{KS9505004}
\begin{equation}\label{eq:KS-CW}
  \frac{3\Nc}{k+1} < \Nf < \frac{3\Nc}{k}\,.
\end{equation}
For $k = 1$, this reduces to the Seiberg window
$\frac{3}{2}\Nc < \Nf < 3\Nc$.

\subsection{Generation bound as a function of $k$}

We identify $\SU{3}_c$ as the magnetic gauge group for the QCD sector:
\begin{equation}
  k\Nf - \Nc = 3
  \qquad\Longrightarrow\qquad
  \Nc = k\Nf - 3 = 2k\ngen - 3\,,
\end{equation}
using $\Nf = 2\ngen$.

The electric gauge group is
\begin{equation}\label{eq:KS-electric}
  \SU{2k\ngen - 3}\,.
\end{equation}

\paragraph{Non-abelian condition.}
$2k\ngen - 3 \geq 2$ requires $\ngen \geq \frac{5}{2k}$:
\begin{equation}\label{eq:KS-bound}
  \ngen \geq \left\lceil \frac{5}{2k} \right\rceil
  \;=\; \begin{cases}
    3 & k = 1\,, \\
    2 & k = 2\,, \\
    1 & k \geq 3\,.
  \end{cases}
\end{equation}

\paragraph{Conformal window.}
Substituting $\Nc = 2k\ngen - 3$ into~\eqref{eq:KS-CW}:
\begin{align}
  \text{Lower:}&\quad 2\ngen > \frac{3(2k\ngen - 3)}{k+1}
  = \frac{6k\ngen - 9}{k+1}\,,
  \label{eq:KS-lower} \\
  \text{Upper:}&\quad 2\ngen < \frac{3(2k\ngen - 3)}{k}
  = 6\ngen - \frac{9}{k}\,.
  \label{eq:KS-upper}
\end{align}

\textbf{Lower bound.}  From~\eqref{eq:KS-lower}:
$2\ngen(k+1) > 6k\ngen - 9$, so $(2 - 4k)\ngen > -9$.  For $k \geq 1$:
$2 - 4k < 0$, giving
\begin{equation}\label{eq:KS-lower-ng}
  \ngen < \frac{9}{4k - 2}\,.
\end{equation}

\textbf{Upper bound.}  From~\eqref{eq:KS-upper}:
$2\ngen < 6\ngen - 9/k$, so $4\ngen > 9/k$, giving
\begin{equation}\label{eq:KS-upper-ng}
  \ngen > \frac{9}{4k}\,.
\end{equation}

The conformal window in the $\ngen$ variable is
\begin{equation}\label{eq:KS-CW-ng}
  \frac{9}{4k} < \ngen < \frac{9}{4k - 2}\,.
\end{equation}

\subsection{Survey over $k$}

\begin{center}
\renewcommand{\arraystretch}{1.5}
\begin{tabular}{c c c c c}
\toprule
$k$ & Non-abelian bound & CW range for $\ngen$ &
  Integer $\ngen$ in CW & Status \\
\midrule
1 & $\ngen \geq 3$ & $2.25 < \ngen < 4.5$ & $\{3, 4\}$ & viable \\
2 & $\ngen \geq 2$ & $1.125 < \ngen < 1.5$ & $\emptyset$ &
  \textbf{ruled out} \\
3 & $\ngen \geq 1$ & $0.75 < \ngen < 0.9$ & $\emptyset$ &
  \textbf{ruled out} \\
$k \geq 4$ & $\ngen \geq 1$ & $< 1$ & $\emptyset$ &
  \textbf{ruled out} \\
\bottomrule
\end{tabular}
\end{center}

For $k \geq 2$, the conformal window is too narrow to contain any
positive integer.  The width of the window in $\ngen$ is
\begin{equation}
  \Delta\ngen = \frac{9}{4k-2} - \frac{9}{4k}
  = \frac{9 \cdot 2}{4k(4k-2)} = \frac{9}{k(4k-2)}\,,
\end{equation}
which decreases rapidly: $\Delta\ngen(k{=}1) = 2.25$,
$\Delta\ngen(k{=}2) = 0.375$, $\Delta\ngen(k{=}3) \approx 0.15$.
Only $k = 1$ has a window wide enough to contain integers.

\begin{remark}[Physical interpretation]
The narrowing of the conformal window for $k \geq 2$ is a consequence
of the adjoint superpotential $W = \Tr(\Phi^{k+1})$, which constrains
the anomalous dimension of~$\Phi$ and thereby restricts the range of
$\Nf$ compatible with the IR fixed point.  Higher~$k$ corresponds to a
more relevant deformation that squeezes the window from both sides.
The result is that the only Kutasov--Schwimmer theory compatible with
integer generation numbers is $k = 1$, which reduces to ordinary
Seiberg duality.
\end{remark}


%% ================================================================
%%  SECTION 7: MASTER TABLE AND INTERSECTION
%% ================================================================
\section{Master Comparison Table}
\label{sec:master-table}

We collect the generation bounds from all duality types in a single
table.

\begin{table}[h!]
\centering
\renewcommand{\arraystretch}{1.6}
\begin{tabular}{@{}l l l c c l@{}}
\toprule
\textbf{Duality} & \textbf{Magnetic} & \textbf{Electric} &
  \textbf{Group thy.} & \textbf{CW} &
  \textbf{Combined} \\
& \textbf{= SM sector} & \textbf{gauge group} &
  \textbf{bound} & \textbf{range} & \\
\midrule
SU Seiberg & $\SU{3}_c$ & $\SU{2\ngen - 3}$ &
  $\ngen \geq 3$ & $\{3,4\}$ & $\ngen \in \{3,4\}$ \\
SU Pati--Salam & $\SU{3}_c$ via $\SU{4}$ & $\SU{2\ngen - 4}$ &
  $\ngen \geq 3$ & $\{3,4,5,6\}$ & $\ngen \in \{3,...,6\}$ \\
$\Sp{}$ (IS) & $\Sp{1} = \SU{2}_L$ & $\Sp{2\ngen - 3}$ &
  $\ngen \geq 2$ & $\{2, 3\}$ & $\ngen \in \{2,3\}$ \\
$\SO{}$ (IS) & $\SO{6} \cong \SU{4}$ & $\SO{2\ngen - 2}$ &
  $\ngen \geq 3$ & $\{4,5\}$ & $\ngen \in \{4,5\}$ \\
KS ($k=1$) & $\SU{3}_c$ & $\SU{2\ngen - 3}$ + adj &
  $\ngen \geq 3$ & $\{3,4\}$ & $\ngen \in \{3,4\}$ \\
KS ($k \geq 2$) & $\SU{3}_c$ & $\SU{2k\ngen - 3}$ + adj &
  various & $\emptyset$ & \textbf{ruled out} \\
\bottomrule
\end{tabular}
\caption{Generation bounds across all Seiberg-type duality types.
``Group thy.''\ is the lower bound from requiring a non-abelian electric
gauge group.  ``CW'' is the range of integer $\ngen$ compatible with the
conformal window (including boundary cases).  ``Combined'' is the
intersection.  The Pati--Salam CW includes the $\ngen = 3$ boundary case
following~\cite{Sannino2011}.}
\label{tab:master}
\end{table}

\subsection{Intersection of all bounds}

The viable duality types for the Standard Model are those that can
accommodate $\SU{}$ or $\Sp{}$ gauge groups (since the SM gauge groups
are $\SU{3}_c$, $\SU{2}_L \cong \Sp{1}$, and $\UU{1}_Y$).  The SO
route is excluded because it requires $\ngen \geq 4$.  KS with $k \geq 2$
is excluded because it has no integer solutions.

From the remaining viable routes:
\begin{equation}\label{eq:intersection}
  \underbrace{\ngen \geq 3}_{\text{SU route}}
  \;\;\cap\;\;
  \underbrace{2 \leq \ngen \leq 3}_{\text{Sp route}}
  \qquad\Longrightarrow\qquad
  \boxed{\ngen = 3}\,.
\end{equation}

This is the central result of this paper: \textbf{if both the QCD and
electroweak sectors of the Standard Model admit UV completions via
Seiberg-type dualities, the number of generations is uniquely determined
to be three.}

\subsection{Role of the SO exclusion}

The SO~route provides a complementary constraint: it demonstrates that
\emph{not all} duality types are compatible with the Standard Model
spectrum.  Specifically, the $+4$ shift in the SO dual gauge group
formula~\eqref{eq:SO-dual} pushes the conformal window to $\ngen \geq 4$,
making $\SO{}$ duality incompatible with $\ngen = 3$.

This exclusion is physically meaningful: it implies that the duality
frame of the Standard Model, \emph{if one exists}, must be of the
$\SU{}$ or $\Sp{}$ type---precisely the types that match the Standard
Model gauge groups $\SU{3}_c$ and $\SU{2}_L \cong \Sp{1}$.  The
orthogonal duality is a ``wrong fit'' for the Standard Model.

\subsection{Robustness of the $\ngen = 3$ selection}

The result $\ngen = 3$ depends on:
\begin{enumerate}
  \item The SU~route lower bound $\ngen \geq 3$ (from requiring
    $\SU{2\ngen - 3}$ or $\SU{2\ngen - 4}$ to be non-abelian).
  \item The Sp~route upper bound $\ngen \leq 3$ (from the conformal
    window of $\Sp{2\ngen - 3}$).
\end{enumerate}
The lower bound is purely group-theoretic and does not depend on
dynamical assumptions.  The upper bound depends on the SUSY conformal
window formula~\eqref{eq:Sp-CW}, which is exact in the SUSY context
but whose extrapolation to non-SUSY theories is uncertain.

If the non-SUSY conformal window for Sp theories is wider than the
SUSY one (as suggested by some lattice studies), the upper bound on
$\ngen$ would relax.  However, the lower bound $\ngen \geq 3$ from the SU
route is robust, and the lower bound $\ngen \geq 2$ from the Sp route is
purely group-theoretic.


%% ================================================================
%%  SECTION 8: CONFORMAL WINDOW DETAILS
%% ================================================================
\section{Conformal Window Compatibility}
\label{sec:conformal-details}

For completeness, we present the conformal window analysis for each
duality type at the physically relevant values of~$\ngen$.

\begin{table}[h!]
\centering
\renewcommand{\arraystretch}{1.4}
\begin{tabular}{@{}l c c c c c c@{}}
\toprule
\textbf{Route} & $\ngen$ & $\Nc$ & $\Nf$ &
  $b_0^{\text{el}}$ & CW lower & CW upper \\
\midrule
\multirow{3}{*}{SU Seiberg} & 3 & 3 & 6 & 7 & 4.5 & 9 \\
  & 4 & 5 & 8 & 9.67 & 7.5 & 15 \\
  & 5 & 7 & 10 & 12.33 & 10.5 & 21 \\
\midrule
\multirow{3}{*}{SU PS} & 3 & 2 & 6 & 0 & 3 & 6 \\
  & 4 & 4 & 8 & 4 & 6 & 12 \\
  & 5 & 6 & 10 & 8 & 9 & 18 \\
\midrule
\multirow{3}{*}{Sp} & 2 & 1 & 4 & 2 & 3 & 6 \\
  & 3 & 3 & 6 & 6 & 6 & 12 \\
  & 4 & 5 & 8 & 10 & 9 & 18 \\
\midrule
\multirow{3}{*}{SO} & 3 & 4 & 6 & 0 & 3 & 6 \\
  & 4 & 6 & 8 & 4 & 6 & 12 \\
  & 5 & 8 & 10 & 8 & 9 & 18 \\
\bottomrule
\end{tabular}
\caption{Conformal window analysis for each duality route.
$b_0^{\text{el}}$ is the one-loop beta function coefficient of the
electric theory.  ``CW lower'' and ``CW upper'' are the boundaries
$\frac{3}{2} \times (\text{dual Coxeter})$ and
$3 \times (\text{dual Coxeter})$ expressed in terms of $\Nf$.
The entry is ``inside CW'' when
lower~${}<\Nf<{}$~upper, ``at boundary'' when $\Nf$ equals a
boundary, and ``outside'' otherwise.  For the SU PS route,
$b_0 = 2\Nc - \Nf$ (with adjoint).}
\label{tab:conformal}
\end{table}

\textbf{Key observations:}
\begin{itemize}
  \item \textbf{SU Seiberg, $\ngen = 3$:} $\Nf = 6$ is well inside the
    conformal window ($4.5 < 6 < 9$).  This is the most favourable case.
  \item \textbf{Sp, $\ngen = 3$:} $\Nf = 6$ equals the lower CW boundary
    exactly (free magnetic phase).
  \item \textbf{SO, $\ngen = 3$:} $b_0 = 0$; the electric theory is not
    asymptotically free.  The theory is at the upper boundary of the
    CW, which is the loss-of-AF boundary.  This confirms the SO
    exclusion.
  \item \textbf{SU PS, $\ngen = 3$:} $b_0 = 0$; similar to SO, at the
    boundary of AF.  The difference from the SU Seiberg route is that
    the PS route uses a smaller electric gauge group ($\SU{2}$ vs
    $\SU{3}$) due to the Pati--Salam embedding's additional offset.
\end{itemize}


%% ================================================================
%%  SECTION 9: DISCUSSION
%% ================================================================
\section{Discussion}
\label{sec:discussion}

\subsection{The uniqueness of $\ngen = 3$}

Our analysis demonstrates that $\ngen = 3$ is the unique number of
generations compatible with the intersection of Seiberg-type duality
bounds across the $\SU{}$ and $\Sp{}$ routes.  The result can be
understood as follows:
\begin{itemize}
  \item The \emph{SU~route} provides a lower bound ($\ngen \geq 3$) because
    too few generations would make the electric gauge group abelian or
    trivial.
  \item The \emph{Sp~route} provides an upper bound ($\ngen \leq 3$)
    because too many generations would push the electric Sp theory below
    the conformal window, where the duality ceases to be operative.
  \item These two constraints ``squeeze'' from opposite sides, leaving
    $\ngen = 3$ as the only integer solution.
\end{itemize}

The SO~route provides independent confirmation of this picture: it
requires $\ngen \geq 4$, which is incompatible with $\ngen = 3$ and
therefore \emph{self-consistently excludes SO} as a viable duality frame
for the Standard Model.

\subsection{Model dependence}

The generation bounds depend on how the Standard Model matter content
is embedded in the duality framework:
\begin{enumerate}
  \item \textbf{Flavour counting:} The number of ``flavours'' $\Nf$
    depends on whether we count quarks only (direct QCD route) or
    quarks and leptons together (Pati--Salam route).  Both choices give
    $\ngen \geq 3$ for the SU route.
  \item \textbf{Embedding of $\SU{2}_L$:} The Sp~route requires
    identifying all $\SU{2}_L$ doublets (quarks and leptons) as
    fundamentals of the electric Sp group, which is unambiguous.
  \item \textbf{Extension to non-SUSY:} The conformal window bounds are
    exact in the SUSY context.  Their non-SUSY extrapolation may modify
    the quantitative bounds but is unlikely to change the qualitative
    picture of the SU lower bound vs.\ Sp upper bound squeeze.
\end{enumerate}

\subsection{Comparison with other generation-number arguments}

Previous arguments for $\ngen = 3$ include:
\begin{itemize}
  \item \textbf{Anomaly cancellation:} The SM is anomaly-free for any
    $\ngen$; this does not constrain the generation number.
  \item \textbf{CP violation:} The CKM matrix has a physical phase only
    for $\ngen \geq 3$, but this does not explain \emph{why} the phase
    should be nonzero.
  \item \textbf{Asymptotic freedom of QCD:} $b_0 = 11 - 2\Nf/3 > 0$
    requires $\Nf < 16.5$, \ie, $\ngen \leq 8$.  This is an upper bound
    but a weak one.
  \item \textbf{Cosmological and precision constraints:}
    Limits from $Z$-boson width, BBN, and CMB give $\ngen = 3$ for
    light neutrinos, but these are experimental measurements, not
    theoretical explanations.
\end{itemize}

The duality argument is qualitatively different: it derives $\ngen = 3$
from the \emph{internal consistency} of identifying the SM as a magnetic
dual, without reference to specific experimental data.  However, it
rests on the unproven assumption that Seiberg-type dualities extend to
non-supersymmetric theories.

\subsection{Testability}

The duality framework makes additional predictions beyond the generation
number:
\begin{itemize}
  \item \textbf{Composite Higgs:} In the SU~route, the Higgs is a
    composite meson with specific quantum numbers and coupling
    structure~\cite{CacciapagliaSannino2024}.
  \item \textbf{Scalar diquarks and leptoquarks:} In the Sp~route, the
    meson content includes colour-triplet scalars that could be
    observable at colliders~\cite{Paper1}.
  \item \textbf{Conformal window boundary:} The electric theory for
    $\ngen = 3$ sits at or near the conformal window boundary for both
    the SU and Sp routes.  Lattice studies of theories in this
    regime could test the existence of the IR fixed point.
\end{itemize}


%% ================================================================
%%  SECTION 10: CONCLUSIONS
%% ================================================================
\section{Conclusions}
\label{sec:conclusions}

We have systematically derived the generation bound $\ngen$ for the
Standard Model across all established $\mathcal{N}=1$ Seiberg-type
dualities.  Our results are:

\begin{enumerate}
  \item \textbf{SU Seiberg duality} ($\SU{3}_c$ as magnetic):
    electric $\SU{2\ngen - 3}$, bound $\ngen \geq 3$, CW-compatible for
    $\ngen \in \{3, 4\}$.

  \item \textbf{SU Pati--Salam} ($\SU{3}_c$ via $\SU{4}_{\mathrm{PS}}$):
    electric $\SU{2\ngen - 4}$, bound $\ngen \geq 3$, CW-compatible for
    $\ngen \in \{3, 4, 5, 6\}$.

  \item \textbf{Sp duality} ($\SU{2}_L \cong \Sp{1}$ as magnetic):
    electric $\Sp{2\ngen - 3}$, bound $\ngen \geq 2$, CW-compatible for
    $\ngen \in \{2, 3\}$.  Upper bound $\ngen \leq 3$.

  \item \textbf{SO duality} ($\SO{6} \cong \SU{4}$ as magnetic):
    electric $\SO{2\ngen - 2}$, bound $\ngen \geq 3$, CW-compatible only
    for $\ngen \in \{4, 5\}$.  \textbf{Excludes $\ngen = 3$.}

  \item \textbf{Kutasov--Schwimmer} ($k \geq 2$): no integer $\ngen$ in
    the conformal window.  \textbf{Ruled out.}  Only $k = 1$ (Seiberg
    duality) is viable.
\end{enumerate}

Within the chosen duality-embedding framework, the intersection of the SU and Sp bounds uniquely selects
\begin{equation}
  \boxed{\ngen = 3}\,.
\end{equation}
The SO~route independently confirms the consistency of this result by
excluding itself as a viable duality frame for $\ngen = 3$.

If the conjectured non-supersymmetric extension of Seiberg duality holds,
the three-generation structure of the Standard Model is not a free
parameter but a \emph{mathematical consequence} of requiring both the QCD
and electroweak sectors to admit self-consistent UV completions via
electric--magnetic duality.


%% ================================================================
%%  ACKNOWLEDGEMENTS
%% ================================================================
\section*{Acknowledgements}

This work is part of the Dualised Standard Model lecture programme.
The analysis builds on the Cacciapaglia--Sannino
construction~\cite{CacciapagliaSannino2024} and the Sp-route analysis
of~\cite{Paper1}.


%% ================================================================
%%  REFERENCES
%% ================================================================
\begin{thebibliography}{99}

\bibitem{Seiberg1995}
N.~Seiberg,
``Electric-magnetic duality in supersymmetric non-Abelian gauge theories,''
\textit{Nucl.\ Phys.\ B} \textbf{435} (1995) 129--146,
\href{https://arxiv.org/abs/hep-th/9411149}{hep-th/9411149}.

\bibitem{IS9506101}
K.~Intriligator and N.~Seiberg,
``Duality, monopoles, dyons, confinement and oblique confinement
in supersymmetric $\SO{\Nc}$ gauge theories,''
\textit{Nucl.\ Phys.\ B} \textbf{444} (1995) 125--160,
\href{https://arxiv.org/abs/hep-th/9503179}{hep-th/9503179}.

\bibitem{IS1995}
K.~Intriligator and N.~Seiberg,
``Lectures on supersymmetric gauge theories and electric-magnetic duality,''
\textit{Nucl.\ Phys.\ Proc.\ Suppl.} \textbf{45BC} (1996) 1--28,
\href{https://arxiv.org/abs/hep-th/9509066}{hep-th/9509066}.

\bibitem{KS9505004}
D.~Kutasov and A.~Schwimmer,
``On duality in supersymmetric Yang-Mills theory,''
\textit{Phys.\ Lett.\ B} \textbf{354} (1995) 315--321,
\href{https://arxiv.org/abs/hep-th/9505004}{hep-th/9505004}.

\bibitem{KSS9507040}
D.~Kutasov, A.~Schwimmer, and N.~Seiberg,
``Chiral rings, singularity theory and electric-magnetic duality,''
\textit{Nucl.\ Phys.\ B} \textbf{459} (1996) 455--496,
\href{https://arxiv.org/abs/hep-th/9510222}{hep-th/9510222}.

\bibitem{CacciapagliaSannino2024}
G.~Cacciapaglia and F.~Sannino,
``The Standard Model as a Seiberg dual,''
\href{https://arxiv.org/abs/2407.17281}{arXiv:2407.17281} (2024).

\bibitem{Sannino2011}
F.~Sannino,
``Magnetic S-parameter,''
\textit{Phys.\ Rev.\ Lett.} \textbf{105} (2010) 232002,
\href{https://arxiv.org/abs/1007.0254}{arXiv:1007.0254}.

\bibitem{PS1974}
J.~C.~Pati and A.~Salam,
``Lepton number as the fourth 'color',''
\textit{Phys.\ Rev.\ D} \textbf{10} (1974) 275--289.

\bibitem{AMPS2011}
A.~Armoni, M.~Shifman, and G.~Veneziano,
``Refining the proof of planar equivalence,''
\textit{Phys.\ Rev.\ D} \textbf{71} (2005) 045015,
\href{https://arxiv.org/abs/hep-th/0412203}{hep-th/0412203}.

\bibitem{Paper1}
``Sp(1) Seiberg duality as an alternative UV completion of the
electroweak sector,''
Notes from the Dualised Standard Model Programme, 2026.

\bibitem{Frampton1992}
P.~H.~Frampton,
``Chiral Duality and the Generation Problem,''
\textit{Phys.\ Rev.\ Lett.} \textbf{69} (1992) 2889.

\end{thebibliography}

\end{document}
